% =====================================================================
\section{Mở đầu}\label{sec:introduction}
% =====================================================================

\subsection{Lý do chọn đề tài}\label{subsec:rationale}

Trong kỷ nguyên dữ liệu lớn, kích thước các ma trận biểu diễn dữ liệu thực tế --
ảnh số, ma trận đánh giá người dùng--sản phẩm, ma trận dữ liệu cảm biến, văn bản
biểu diễn dưới dạng tần số từ -- thường vượt xa khả năng lưu trữ và xử lý trực
tiếp. Trong bối cảnh đó, phân tích giá trị kỳ dị (\emph{Singular Value
Decomposition} -- SVD) là một trong những công cụ trung tâm của đại số tuyến
tính số hiện đại: nó cho phép tách một ma trận bất kỳ thành tích của ba ma trận
có cấu trúc đẹp, từ đó cung cấp cơ sở toán học cho một loạt ứng dụng quan
trọng như nén dữ liệu, khử nhiễu, giảm chiều, hệ thống gợi ý và xử lý ảnh.

Tuy nhiên, thuật toán SVD cổ điển dựa trên giai đoạn tridiagonalize Golub--Kahan
hoặc các biến thể Jacobi có độ phức tạp thời gian
$\mathcal{O}(mn\min(m,n))$ và đòi hỏi truy xuất toàn bộ ma trận nhiều lần.
Khi kích thước ma trận $m, n$ đạt mức $10^5$ trở lên -- vốn rất phổ biến trong
các bộ dữ liệu thực tế như MovieLens hay các ma trận hình ảnh độ phân giải cao
-- cách tiếp cận này trở nên bất khả thi cả về thời gian lẫn bộ nhớ. Năm 2011,
N.~Halko, P.-G.~Martinsson và J.~A.~Tropp đã hệ thống hóa một họ phương pháp
mang tính đột phá có tên \emph{Randomized SVD}: dùng phép chiếu ngẫu nhiên để
rút trích nhanh một không gian con xấp xỉ tốt cho $k$ giá trị kỳ dị lớn nhất,
sau đó chỉ chạy SVD cổ điển trên một ma trận nhỏ hơn nhiều. Phương pháp này có
cơ sở toán học vững chắc, sai số kiểm soát được dưới dạng kỳ vọng và độ lệch
chuẩn, dễ song song hóa và đặc biệt hiệu quả với ma trận thưa.

Đề tài tập trung vào Randomized SVD vì ba lý do căn bản:

\begin{itemize}
    \item \textbf{Tầm quan trọng toán học:} Randomized SVD là một kết quả tiêu
    biểu của đại số tuyến tính số kết hợp với lý thuyết xác suất hiện đại
    (xác suất tập trung, bất đẳng thức ma trận nhân tử). Nó cũng là cầu nối
    quan trọng giữa lý thuyết phổ và các phương pháp Monte~Carlo.
    \item \textbf{Lĩnh vực áp dụng rộng:} học máy (Latent Semantic Analysis,
    Principal Component Analysis), xử lý ảnh (image compression, khử nhiễu),
    hệ thống gợi ý (Netflix Prize), nén tín hiệu, mô hình hóa ngôn ngữ và tin
    sinh học (gene expression analysis).
    \item \textbf{Vấn đề thực tế cụ thể:} bài toán giảm dung lượng lưu trữ và
    rút gọn chiều biểu diễn cho dữ liệu lớn, sao cho vẫn giữ được một tỉ lệ
    năng lượng (information content) $\eta$ cho trước. Đây là một bài toán
    điển hình mà sinh viên kỹ thuật cần thuần thục: cân bằng giữa độ chính
    xác xấp xỉ và chi phí lưu trữ/tính toán.
\end{itemize}

\subsection{Mục tiêu của đề tài}\label{subsec:goals}

Báo cáo hướng tới các kết quả cụ thể, có thể kiểm chứng được sau đây:

\begin{enumerate}
    \item \textbf{Cơ sở lý thuyết:} Trình bày đầy đủ và nghiêm túc cơ sở toán
    học của SVD đầy đủ, Compact SVD và Randomized SVD; nêu rõ ý nghĩa
    hình học cũng như đại số của các đại lượng $U$, $\Sigma$, $V$. Phần này
    được trình bày trong \cref{sec:theory}.

    \item \textbf{Tiêu chuẩn chọn $k$ theo tỉ lệ năng lượng:} Xây dựng và
    biện luận công thức
    \[
        \frac{\sum_{i=1}^{k}\sigma_i^2}{\sum_{i=1}^{r}\sigma_i^2} \;\ge\; \eta,
    \]
    đồng thời phân tích ý nghĩa của tham số $\eta$ trong việc kiểm soát đánh
    đổi giữa độ chính xác và mức độ nén.

    \item \textbf{Cài đặt thuật toán Randomized SVD} (\cref{sec:algorithm}):
    nhóm tự viết phiên bản đầy đủ với các bước phép chiếu ngẫu nhiên Gauss,
    \emph{power iteration} tăng độ tách phổ, trực giao hóa QR và SVD đầy đủ
    trên ma trận nhỏ. \emph{Không} sử dụng trực tiếp các hàm SVD đầy đủ trên
    ma trận gốc.

    \item \textbf{Áp dụng cho ma trận thực kích thước lớn:} sử dụng bộ dữ liệu
    \mlSource{} với $m \times n = \mlM \times \mlN$ và $\mathrm{nnz} = \mlNNZ$
    (vượt yêu cầu $10^5$ phần tử của đề bài), cùng một ảnh xám
    $512\times 512$ làm trường hợp dữ liệu \emph{dày}.

    \item \textbf{Đánh giá định lượng:} đo sai số xấp xỉ
    $\|A - A_k\|_F$ theo chuẩn Frobenius, tỉ lệ năng lượng giữ lại $S_k/\|A\|_F^2$,
    thời gian chạy so với SVD cổ điển và \texttt{scikit-learn}.

    \item \textbf{Phát hiện điều kiện break-even cho nén:} rút ra điều kiện
    định lượng $k_{be} = \mathrm{nnz}(A)/(m+n+1)$ cho biết với loại ma trận
    nào thì Randomized SVD nén được dung lượng -- một kết quả thực tế quan
    trọng nhưng thường bị các tài liệu chuẩn bỏ qua.
\end{enumerate}

\subsection{Phạm vi nghiên cứu}\label{subsec:scope}

\begin{itemize}
    \item \textbf{Đối tượng nghiên cứu:} ma trận thực
    $A \in \mathbb{R}^{m\times n}$ với $m \cdot n \ge 10^5$. Hai loại
    dữ liệu được khảo sát: ma trận \emph{thưa} (rating matrix MovieLens 25M)
    và ma trận \emph{dày} (ảnh xám tiêu chuẩn \texttt{skimage.data.camera}).

    \item \textbf{Mức độ lý thuyết--ứng dụng:} báo cáo trình bày cân bằng giữa
    lý thuyết và ứng dụng: chứng minh chặt chẽ các định lý nền tảng (định lý
    phổ, định lý Eckart--Young, hệ quả về chuẩn Frobenius) đi kèm các
    chứng minh hoặc phác thảo chứng minh; mọi kết quả số đều dựa trên thực
    nghiệm độc lập bằng cài đặt tự viết.

    \item \textbf{Nội dung không xét đến:} các phiên bản phân tán
    (distributed) hay GPU-accelerated của Randomized SVD; các biến thể khối
    (block Lanczos, randomized block Krylov) phức tạp; phân tích sai số xác
    suất ở mức kỹ thuật (chỉ trình bày bound kỳ vọng dạng định tính).
\end{itemize}

Sau \cref{sec:introduction}, báo cáo được tổ chức như sau: \cref{sec:theory}
trình bày cơ sở lý thuyết của SVD, Compact SVD và định lý Eckart--Young;
\cref{sec:algorithm} mô tả chi tiết thuật toán Randomized SVD cùng phân tích
độ phức tạp và sai số; \cref{sec:experiments} báo cáo kết quả thực nghiệm
trên hai bộ dữ liệu thực; cuối cùng \cref{sec:conclusion} tổng kết và đề xuất
hướng phát triển.
